% !TeX root = ../初等数学笔记.tex
% 由主文档通过 \input 载入；请编译主文档。
\chapter*{绪论：研究对象、课程总纲与学习路线}
\addcontentsline{toc}{chapter}{绪论：研究对象、课程总纲与学习路线}

\section*{初等数学研究什么}
\emph{初等数学}（elementary mathematics）以数、式、方程、函数和几何图形为基本对象，研究数量关系、空间关系以及有限过程中的规律。它既包含算术、初等代数和欧氏几何，也包含三角学、解析几何、初步组合计数与概率统计。本书将这些内容组织成一条连续的学习路线：从正确表达问题开始，用运算和推理建立关系，再根据关系计算、证明、估计或建模。

例如，“一定周长的长方形怎样使面积最大”可以用代数式、二次函数或不等式处理；“怎样测量无法直接到达的两点间距离”可以用相似三角形、三角函数或坐标处理。同一个问题的不同解法揭示不同结构，学习目标因而不仅是得到答案，还包括说明变量范围、变形依据、结论条件与方法之间的联系。

\section*{全书定位与当前版本}
本书面向希望系统补足初等数学基础、并准备学习大学数学的读者。先修要求是整数和分数的基本运算、常用度量单位及对简单几何图形的直观认识；不预设微积分、线性代数或抽象代数。基础内容覆盖中学数学的主要知识，并以本科生需要的证明意识、结构理解和建模表达加以统一；较深的凸性、图论与极限讨论可作为选读。

本书按本科阶段基础巩固与知识衔接的需要编写，正文覆盖全部二十四章，以定义、主要结论、关键证明和代表例题建立主线，章末配置习题与答案要点。论述尽量紧凑，但保留变量范围、等价变形依据和取等条件。实数系的基本结构、欧氏几何公理及部分几何基础定理作为先修事实使用；实指数延拓、连续函数介值性质和球面面积等未完整证明的结论在使用处明确说明边界。

\section*{课程总纲}
全书分六篇、二十四章，按照知识依赖安排。以下总纲概括课程结构，随后各篇展开正文。
\begin{description}
  \item[第一篇：数学语言、数与式]
  第\ref{chap:em-language}章，集合、逻辑与证明；第\ref{chap:em-numbers}章，数系、运算与近似；第\ref{chap:em-integers}章，整除、素数与同余初步；第\ref{chap:em-expressions}章，代数式与多项式。建立表达、计算和论证的共同基础。
  \item[第二篇：方程、不等式与函数]
  第\ref{chap:em-equations}章，方程与方程组；第\ref{chap:em-solving-inequalities}章，不等式与不等式组；第\ref{chap:em-functions}章，函数与图像；第\ref{chap:em-exponential}章，幂、指数与对数。以解集和定义域贯穿代数变形与函数研究。
  \item[第三篇：数列、估计与有限方法]
  第\ref{chap:em-sequences}章，数列、递推与求和；第\ref{chap:em-inequality-methods}章，初等不等式与最值；第\ref{chap:em-counting}章，计数原理与二项式定理；第\ref{chap:em-discrete}章，有限问题中的证明方法。联系规律发现、严格证明与有效计算。
  \item[第四篇：平面与空间几何]
  第\ref{chap:em-triangles}章，直线、三角形与相似；第\ref{chap:em-circles}章，四边形与圆；第\ref{chap:em-transformations}章，几何变换、轨迹与作图；第\ref{chap:em-solid}章，立体几何与度量。以欧氏几何中的位置关系、全等、相似和度量为主线。
  \item[第五篇：三角、坐标与复数]
  第\ref{chap:em-trigonometry}章，角与三角函数；第\ref{chap:em-triangle-solving}章，三角恒等变换、方程与解三角形；第\ref{chap:em-vectors}章，平面向量、直线与圆的坐标方法；第\ref{chap:em-conics}章，圆锥曲线；第\ref{chap:em-complex}章，复数及其几何意义。用代数语言统一描述长度、方向、位置与旋转。
  \item[第六篇：随机、数据与综合应用]
  第\ref{chap:em-probability}章，有限概率模型；第\ref{chap:em-statistics}章，统计描述与数据解释；第\ref{chap:em-modeling}章，数学建模、综合方法与后续学习。把确定性工具用于带有随机性或现实约束的问题。
\end{description}
附录\ref{app:em-notation}整理符号、公式与方法索引，附录\ref{app:em-study}提供分层自测及专题阅读路径。附录用于定位正文中的定义、条件与推导，不能替代正文学习。

\section*{学习路径与章节依赖}
基础学习沿前两篇顺序推进，随后可分两条线阅读：代数线为数列、不等式、计数，再到概率与统计；几何线为三角形、圆、几何变换、立体几何，再到三角与坐标。第五篇使用第二篇的函数语言及第四篇的几何基础，复数的三角形式还依赖三角恒等式。第\ref{chap:em-discrete}章的提高题、第\ref{chap:em-inequality-methods}章的凸性、第\ref{chap:em-conics}章的综合轨迹题可以在第一遍学习时暂缓。

已有专题笔记中的均值、伯努利不等式与 Jensen 不等式，分别归入第\ref{chap:em-inequality-methods}章的主干或选读内容；裂项求和与双阶乘分别归入第\ref{chap:em-sequences}章和第\ref{chap:em-counting}章；三角公式与解三角形对应第\ref{chap:em-trigonometry}、\ref{chap:em-triangle-solving}章；圆和球的度量对应第\ref{chap:em-circles}、\ref{chap:em-solid}章。正文按本书的先修顺序组织必要内容，较深的专题论证作为补充阅读，不要求读者预先读过这些独立笔记。

\section*{符号、论证与练习约定}
记自然数集为 $\mathbb N=\{0,1,2,\ldots\}$，正整数集为 $\mathbb N_{>0}$，整数、有理数、实数与复数集分别为 $\mathbb Z,\mathbb Q,\mathbb R,\mathbb C$；复数在第\ref{chap:em-complex}章正式引入。未知数通常记为 $x,y$，数列项记为 $a_n$，函数记为 $f$。同一字母在不同章节可能承担不同角色，使用时重新说明；几何中的 $A,B,C$ 通常表示点，三角形内角则写为 $\angle A,\angle B,\angle C$，需要简记时作明确约定。

等号表示相等，$\Longleftrightarrow$ 表示等价，$\Longrightarrow$ 表示单向推出，$\approx$ 表示带有上下文精度要求的近似。方程和不等式的变形始终保留定义域；除以可能为零的量必须分类，平方或开方必须检查符号。几何图形用于理解关系，不能代替证明；随机试验的样本空间和等可能性必须说明；实际数据与教学自拟数据严格区分。

章末习题结合各章内容安排概念辨析、基本计算、证明与反例或综合应用，附录提供分阶段检查任务。基础题训练正确使用定义，提高题训练选择方法，综合题训练在不同表示之间转换。代表例题给出关键依据与验算，章末习题附答案或提示；数值计算和作图用于探索与核验，不能据有限个例子断言一般规律。

